\documentclass[slideColor]{prosper}
\usepackage{amsmath,amsfonts,amsthm}
%\usepackage{epsfig}
%\include{Qcircuit}

%\usepackage{algorithm} 
%\usepackage{algorithmic}

\newcommand{\E}{\mathbf{E}}
\renewcommand{\R}{\mathbb{R}}
\newcommand{\Z}{\mathbb{Z}}

\newcommand{\cA}{\mathcal A}
\newcommand{\cB}{\mathcal B}
\newcommand{\cH}{\mathcal H}
\newcommand{\C}{\mathbb C}
\newcommand{\TVD}{\mathrm{TVD}}

\renewcommand{\>}{\rangle}
\newcommand{\<}{\langle}   


%%%%%%%%%%%%%%%%%%%%%%%%%%%%%%%%%%%%%%%%%%%%%%%%%%%%%%%%%%%%

\title{Efficient Quantum Algorithms for One-Dimensional Infrastructures}
\author{Pawel M. Wocjan\\[1ex]Department of Electrical Engineering and Computer Science\\ University of Central Florida\\Orlando}
\email{wocjan@eecs.ucf.edu}

\begin{document}

\maketitle
%%%%%%%%%%%%%%%%%%%%%%%%%%%%%%%%%%%%%%%%%%%%%%%%%%%%%%%%%%%%%%%%%%%%%%%

\begin{slide}{Joint work with}
\begin{itemize}
\item Pradeep Sarvepalli (University of British Columbia) 

\bigskip
preprint arXiv:1106.6347
 
\bigskip\bigskip
\item P.W. gratefully acknowledges the support from the NSF grant CCF-0726771 and the NSF CAREER Award CCF-0746600

\bigskip
P.S. was sponsored by grants from CIFAR, MITACS, and NSERC
\end{itemize}
\end{slide}

%%%%%%%%%%%%%%%%%%%%%%%%%%%%%%%%%%%%%%%%%%%%%%%%%%%%%%%%%%%%%%%%%%%%%%%%%

\begin{slide}{Motivation and overview}
\vspace{-0.25cm}
\begin{itemize}
\item infrastructures are group-like structures

\medskip
\item they appear in arithmetic geometry in the study of computational problems related to
number fields and function fields over finite fields 

\medskip
\item arithmetic geometry provides a unified treatment of global fields 

\medskip
\item
the most prominent computational tasks are computing 

\begin{itemize}
\item the circumference of the infrastructure and

\medskip
\item generalized discrete logarithms
\end{itemize}

\medskip
\item
both problems are not know to have efficient classical algorithms
\end{itemize}

\end{slide}

%%%%%%%%%%%%%%%%%%%%%%%%%%%%%%%%%%%%%%%%%%%%%%%%%%%%%%%%%%%%%%%%%%%%%%%%%

\begin{slide}{Discrete Infrastructures}
\begin{itemize}
\item there is a simple and efficient quantum algorithm for solving these problems for infrastructures arising from function fields 

\medskip
\item this is because such infrastructures are discrete and can be embedded into finite abelian groups

\medskip
\item we can simply use the efficient quantum algorithm for the abelian hidden subgroup problem to solve these two problems 
and also many other problems related function fields

\end{itemize}

\end{slide}

%%%%%%%%%%%%%%%%%%%%%%%%%%%%%%%%%%%%%%%%%%%%%%%%%%%%%%%%%%%%%%%%%%%%%%%%%

\begin{slide}{1-dim nondiscrete infrastructures}
\begin{itemize}

\medskip
\item $\Rightarrow$ we focus on non-discrete infrastructures

\medskip
\item
such infrastructures arise from number fields

\medskip
\item in this talk, I only consider one-dimensional infrastructures

\medskip
\item such infrastructures arise from number fields of unit rank
\end{itemize}

\end{slide}

%%%%%%%%%%%%%%%%%%%%%%%%%%%%%%%%%%%%%%%%%%%%%%%%%%%%%%%%%%%%%%%%%%%%%%%%%

\begin{slide}{Improvements}
\vspace{-0.4cm}
\begin{itemize}
\item quadratic number fields give rise to one-dimensional infrastructures $\Rightarrow$ 
our algorithms can be used to solve Pell's equation and the principal ideal problem

\medskip
\item in this sense, they generalize Hallgren's quantum algorithms for these problems

\medskip
\item our improvements and extensions:

\begin{itemize}
\item give a conceptually simpler presentation 

\item find algorithms with lower complexity and constant success probability

\item introduce a technical result for analyzing algorithms based on Fourier sampling

\item present a more efficient method for computing in one-dimensional infrastructures
\end{itemize}
\end{itemize}
\end{slide}

%%%%%%%%%%%%%%%%%%%%%%%%%%%%%%%%%%%%%%%%%%%%%%%%%%%%%%%%%%%%%%%%%%%%%%%%%

\begin{slide}{One-dimensional infrastructure}

\begin{itemize}
\item a one-dimensional infrastructure of circumference $R$ is a pair $(X,d)$ where

\begin{itemize}
\item $X=\{x_0,x_1,\ldots,x_{m-1}\}$ is a finite set and

\medskip
\item $d \, : \, X \rightarrow \R/R\Z$ is injective
\end{itemize}

\medskip
\item the elements of $X$ are ordered such that 
\[
0=d(x_0)<d(x_1)<\ldots<d(x_{m-1})<R
\]
we use representatives in $[0,R)$ to order the $d(x_i)$

\medskip
\item imagine that we place the elements around a circle of circumference $R$ starting at the top most point and moving clockwise
\end{itemize}

\end{slide}

%%%%%%%%%%%%%%%%%%%%%%%%%%%%%%%%%%%%%%%%%%%%%%%%%%%%%%%%%%%%%%%%%%%%%%%%%

\begin{slide}{Baby step}

\begin{itemize}
\vspace{-0.25cm}
\item the baby step $\mathrm{bs} : X \rightarrow X$ is defined as
\[
\mathrm{bs}(x_i) = \left\{
\begin{array}{ll}
x_{i+1} & 0 \le i < m-1 \\
x_0     & i=m-1
\end{array}
\right.
\]
\medskip
it gives you the next element on the circle

\medskip
\item the relative distance function $\Delta_{\mathrm{bs}} : X \rightarrow [0,R)$ is such that
\[
d(\mathrm{bs}(x)) = d(x) + \Delta_{\mathrm{bs}}(x)
\]
holds for all $x\in X$

\medskip
it tells you how far you have to travel along the circle to get to the next element
\end{itemize}

\end{slide}

%%%%%%%%%%%%%%%%%%%%%%%%%%%%%%%%%%%%%%%%%%%%%%%%%%%%%%%%%%%%%%%%%%%%%%%%%

\begin{slide}{Giant step}
\vspace{-0.25cm}
\begin{itemize}
\item for $x,y\in X$ consider the set
\[
S_{x,y} = \{ s \in [0,R) \, | \, d(x) + d(y) + s \in d(X) \}
\]

\medskip
\item the giant step $\mathrm{gs} : X \times X \rightarrow X$ is defined as
\[
\mathrm{gs}(x,y)=z 
\]

such that
\[
d(z) = d(x) + d(y) + \min S_{x,y}
\]

\medskip
\item the relative distance function $\Delta_{\mathrm{gs}} \, : \, X \times X \rightarrow [0,R)$ is defined as
\[
\Delta_{\mathrm{gs}}(x,y) = \min S_{x,y}
\]
\end{itemize}

\end{slide}

%%%%%%%%%%%%%%%%%%%%%%%%%%%%%%%%%%%%%%%%%%%%%%%%%%%%%%%%%%%%%%%%%%%%%%%%%

\begin{slide}{Computational problems}
\vspace{-0.25cm}
\begin{itemize}
\item estimate the circumference $R$

\medskip
\item given $x\in X$, estimate $d(x)$

\bigskip
\item these problems are computationally hard for classical computers

\medskip
\item the first generalizes the problem of computing the order of a (black-box) cyclic group

\medskip
\item the second generalizes the problem of determining discrete logarithms in a (black-box) cyclic group
\end{itemize}
\end{slide}

%%%%%%%%%%%%%%%%%%%%%%%%%%%%%%%%%%%%%%%%%%%%%%%%%%%%%%%%%%%%%%%%%%%%%%%%%

\begin{slide}{Finite cyclic groups}
\vspace{-0.25cm}
\begin{itemize}
\item suppose $G=\<g\>$ is a finite cyclic group of order $R$

\medskip
\item $\Rightarrow$ we can form an infrastructure out of $G$ as follows:

\begin{itemize}
\item $X=G$

\medskip
\item $d(h) = \log_g(h)$ 

\medskip
\item the baby step corresponds to multiplication by $g$

\medskip
\item the giant step corresponds to the multiplication of two elements in $G$

\medskip
\item $\Delta_{\mathrm{bs}}$ always takes on the value $1$

\medskip 
\item $\Delta_{\mathrm{gs}}$ always takes on the value $0$
\end{itemize}
\end{itemize}
\end{slide}

%%%%%%%%%%%%%%%%%%%%%%%%%%%%%%%%%%%%%%%%%%%%%%%%%%%%%%%%%%%%%%%%%%%%%%%%%

\begin{slide}{Group arithmetic with $f$-reps}

\begin{itemize}
\item a pair $(x,f)\in X\times \R$ is called an $f$-representation if 
\[
0 \le f < \Delta_{\mathrm{bs}}(x)
\]

\medskip
\item denote the set of all $f$-representations by $\mathrm{Rep}(\mathcal{I})$

\medskip
\item $\mathrm{Rep}(\mathcal{I})$ can be endowed with a group structure such that 
\[
\mathrm{Rep}(\mathcal{I}) \cong \R/R\Z
\]
\end{itemize}
\end{slide}

%%%%%%%%%%%%%%%%%%%%%%%%%%%%%%%%%%%%%%%%%%%%%%%%%%%%%%%%%%%%%%%%%%%%%%%%%

\begin{slide}{HSP over $\R$}

\begin{itemize}
\item we have a surjective group homomorphism
\[
h : \R \rightarrow \mathrm{Rep}(\mathcal{I})
\]
with 
\[
\ker h = R \Z
\]

\medskip
\item $\Rightarrow$ computing the circumference corresponds to determining the subgroup of $\R$ hidden by $h$

\medskip
\item this is the intuition behind our algorithm

\medskip
\item but computing $h$ is impossible due to finite precision

\medskip
\item this is because we can only approximate $\Delta_{\mathrm{bs}}$ and $\Delta_{\mathrm{gs}}$
\end{itemize}

\end{slide}

%%%%%%%%%%%%%%%%%%%%%%%%%%%%%%%%%%%%%%%%%%%%%%%%%%%%%%%%%%%%%%%%%%%%%%%%%

\begin{slide}{Computable version $h_N$ of $h$}

\vspace{-0.25cm}
\begin{itemize}
\item Let $h_N : \Z \rightarrow X \times \Z$ be defined as
\[
h_N(i) = (x,\lfloor f N \rfloor)
\]
where $s$ is some random offset and $h(i/N + s) = (x,f)$

\medskip
\item we can compute $h_N(i)$ exactly for $i\in W=\{0,\ldots,q-1\}$ 

\medskip
\item for more than half $(x,k) \in f(W)$, the corresponding preimage has the special form
\[
f^{-1}(x,k) = \{ \lfloor k + j N R \rceil \mid j \in \mathcal{J} \} 
\]

$\lfloor \cdot \rceil$ is a function assigning to each $u\in\R$ either $\lfloor u \rfloor$ or $\lceil u \rceil$

\medskip
\item $\Rightarrow$ evaluating $h_N$ in superposition over $W$ and measuring induces a pseudo-periodic state
\end{itemize}

\end{slide}

%%%%%%%%%%%%%%%%%%%%%%%%%%%%%%%%%%%%%%%%%%%%%%%%%%%%%%%%%%%%%%%%%%%%%%%%%

\begin{slide}{Pseudo-periodic states}
\vspace{-0.25cm}
\begin{itemize}
\item a quantum state in $\C^q$ is called $\epsilon$-pseudo-periodic with period $S$ it is of the following form
\[
|\psi\> = \frac{1}{\sqrt{\mathcal{J}}} \sum_{j\in\mathcal{J}} | \lfloor k + j S\rceil \>
\]
\begin{itemize}
\item $p=\lfloor q/S+1 \rfloor$

\medskip
\item $\mathcal{J}\subseteq\{0,1,\ldots,p-1\}$

\medskip
\item $|\mathcal{J}| \ge \epsilon p$

\medskip
\item $k\in\{0,1,\ldots,\lfloor S \rfloor\}$
\end{itemize}

\medskip
\item we can find a rational number $\hat{S}$ such that $|\hat{S}-S|\le 1$ by Fourier sampling $\epsilon$-pseudo-periodic states
\end{itemize}

\end{slide}

%%%%%%%%%%%%%%%%%%%%%%%%%%%%%%%%%%%%%%%%%%%%%%%%%%%%%%%%%%%%%%%%%%%%%%%%%

\begin{slide}{Fourier sampling}
\vspace{-0.25cm}
\begin{itemize}
\item the lemma ``perturbed geometric sum with missing terms'' is at the heart of the analysis

\begin{itemize}
\item let $\delta\in(-1,1)$

\medskip
\item $\omega$ be the $p$th root of unity $e^{2\pi i/p}$

\medskip
\item $\theta$ an arbitrary real-valued function defined on $\mathcal{J} \subseteq \{0,1,\ldots,p-1\}$ with $|\theta_j|\le \frac{p}{32}$

\item {}
%\vspace{-0.25cm}
\[
|\mathcal{J}|\ge \frac{1-\mathrm{sinc}(\pi\delta)}{1-2\sin(\pi/32)}\, p
\]
\end{itemize}
\[
\!\!\!
\Rightarrow
\frac{1}{|\mathcal{J}|^2} \left| \sum_{j\in\mathcal{J}} \omega^{\delta j + \theta_j} \right|^2 \! \ge
\left(
1 - 2\sin(\pi/32) - \big(1-\mathrm{sinc}(\pi\delta)\big) \frac{p}{|\mathcal{J}|}
\right)^2
\]
\end{itemize}

\end{slide}

%%%%%%%%%%%%%%%%%%%%%%%%%%%%%%%%%%%%%%%%%%%%%%%%%%%%%%%%%%%%%%%%%%%%%%%%%

\begin{slide}{Estimating the period/circumference}

\begin{itemize}
\vspace{-0.4cm}
\item INPUT:

\begin{itemize}
\item an upper bound $M$ on $S>1$ and $q$ an integer such that $M^2\leq q< 2M^2$

\medskip
\item a pair of $\epsilon$-periodic states in $\C^q$ with period $S$
\end{itemize}

\bigskip
\item for each $\epsilon$-pseudo-periodic state,  apply a Fourier transform over $\Z_q$ and measure to obtain $c$ and $d$

\medskip
\item compute the convergents $c_i/d_i$ of $c/d$ where $d_i\leq  \lfloor q/32 \rfloor$

\medskip
\item return $\mathcal{L}= \Big\{ [c_iq/c] \mid d_i\leq \lfloor q/32 \rfloor \Big\}$ as candidates for $S$

\end{itemize}
\end{slide}

%%%%%%%%%%%%%%%%%%%%%%%%%%%%%%%%%%%%%%%%%%%%%%%%%%%%%%%%%%%%%%%%%%%%%%%%%

\begin{slide}{Success probability}

\begin{itemize}
\item our algorithm has constant success probability since $\mathcal{L}$ contains an element $\hat{S}$ with $|\hat{S}-S|<1$ with probability $\Omega(1)$ 

\medskip
\item Hallgren's algorithm succeeds with only $\Omega(1/\log^4(M))$ and works only if $\log M \ge 256$
\end{itemize}
\end{slide}

%%%%%%%%%%%%%%%%%%%%%%%%%%%%%%%%%%%%%%%%%%%%%%%%%%%%%%%%%%%%%%%%%%%%%%%%%

\begin{slide}{Thank you!}

\end{slide}

\end{document} 